\documentclass[a4paper,12pt]{article}
\usepackage[utf8]{inputenc}
\usepackage[T1]{fontenc}
\usepackage{amsmath}
\usepackage{booktabs}
\usepackage{graphicx}
\usepackage{geometry}
\usepackage{xcolor}
\usepackage{listings}
\usepackage{hyperref}
\usepackage{longtable}
\usepackage{array}
\usepackage{float}

\geometry{margin=2.5cm}
\hypersetup{colorlinks=true, linkcolor=blue!50!black, urlcolor=blue!50!black}

\renewcommand{\tablename}{Tabela}
\renewcommand{\lstlistingname}{Código}



% ---------- Configuração do realce de sintaxe (C) ----------

\definecolor{codebg}{RGB}{247,247,247}
\definecolor{codekw}{RGB}{0,0,180}
\definecolor{codecomment}{RGB}{0,128,0}
\definecolor{codestring}{RGB}{163,21,21}
\definecolor{framecolor}{RGB}{200,200,200}

\lstdefinestyle{cstyle}{
    language=C,
    backgroundcolor=\color{codebg},
    basicstyle=\ttfamily\footnotesize,
    keywordstyle=\color{codekw}\bfseries,
    commentstyle=\color{codecomment}\itshape,
    stringstyle=\color{codestring},
    numberstyle=\tiny\color{gray},
    numbers=left,
    numbersep=6pt,
    frame=single,
    rulecolor=\color{framecolor},
    breaklines=true,
    breakatwhitespace=false,
    showstringspaces=false,
    tabsize=4,
    captionpos=b,
    columns=flexible,
    morekeywords={size_t},
    extendedchars=true,
    literate=
        {á}{{\'a}}1 {à}{{\`a}}1 {â}{{\^a}}1 {ã}{{\~a}}1
        {é}{{\'e}}1 {ê}{{\^e}}1 {í}{{\'i}}1
        {ó}{{\'o}}1 {ô}{{\^o}}1 {õ}{{\~o}}1
        {ú}{{\'u}}1 {ü}{{\"u}}1 {ç}{{\c c}}1
        {Á}{{\'A}}1 {À}{{\`A}}1 {Â}{{\^A}}1 {Ã}{{\~A}}1
        {É}{{\'E}}1 {Ê}{{\^E}}1 {Í}{{\'I}}1
        {Ó}{{\'O}}1 {Ô}{{\^O}}1 {Õ}{{\~O}}1
        {Ú}{{\'U}}1 {Ü}{{\"U}}1 {Ç}{{\c C}}1
}



\title{Tarefa 10 -- Mecanismos de Sincronização com OpenMP}
\author{Israel Soares de Castro Filho \\ Matrícula: 20250070780}
\date{}

\begin{document}
\maketitle

\section{Introdução}

Este relatório documenta a Tarefa 10, cujo objetivo é estudar, de forma
prática, os principais mecanismos de sincronização disponíveis no
OpenMP para proteger uma variável compartilhada em regiões paralelas.
Como estudo de caso, mantém-se o problema já utilizado na Tarefa~8:
a estimativa estocástica de $\pi$ pelo método de Monte Carlo, no qual
cada thread sorteia pontos $(x,y)$ no quadrado $[0,1]\times[0,1]$ e
conta quantos caem dentro do quarto de círculo de raio unitário. A
razão entre acertos e total de pontos, multiplicada por 4, converge
para $\pi$.

O foco desta tarefa não é a precisão da estimativa em si, mas sim
\textbf{como} a contagem de acertos é somada entre as threads e o
impacto disso no tempo de execução. Foram implementadas e comparadas
cinco versões do estimador, variando a forma de acumulação do
contador de acertos e o mecanismo de exclusão mútua utilizado.

\section{Metodologia}

A partir do código de referência da Tarefa 8, foram implementadas
5 versões da função de estimativa de $\pi$, todas paralelizadas com
\texttt{\#pragma omp parallel} e \texttt{\#pragma omp for}, variando
apenas a estratégia de acumulação do contador de acertos:

\begin{itemize}
    \item \textbf{V1 -- privado + \texttt{rand()} + critical}: cada
    thread mantém um contador local (\texttt{acertos\_local}) e soma
    seu valor à variável global apenas \textbf{uma vez}, ao final da
    região paralela, dentro de uma seção \texttt{\#pragma omp critical}.
    Utiliza o gerador \texttt{rand()} da biblioteca padrão (não
    thread-safe, herdado da Tarefa 8 como baseline de referência).

    \item \textbf{V2 -- privado + xorshift32 + critical}: idêntico a
    V1 na estratégia de sincronização (contador privado, um único
    \texttt{critical} por thread ao final), mas usando o gerador
    Xorshift32 com estado privado por thread em vez de \texttt{rand()}.

    \item \textbf{V3 -- compartilhado + xorshift32 + critical}: o
    contador de acertos é uma única variável \textbf{compartilhada},
    incrementada dentro de uma seção \texttt{\#pragma omp critical} a
    \textbf{cada acerto} (isto é, potencialmente milhões de entradas na
    seção crítica, e não apenas uma por thread).

    \item \textbf{V4 -- compartilhado + xorshift32 + atomic}: mesma
    estrutura de V3 (contador compartilhado, atualizado a cada acerto),
    trocando a diretiva \texttt{critical} pela diretiva
    \texttt{\#pragma omp atomic}.

    \item \textbf{V5 -- xorshift32 + reduction}: nenhuma diretiva de
    exclusão mútua é escrita manualmente; o acúmulo é delegado à
    cláusula \texttt{reduction(+:total\_acertos)} do próprio
    \texttt{\#pragma omp for}, que cria cópias privadas do contador
    para cada thread e as combina automaticamente ao final do laço.
\end{itemize}

\textbf{Observação sobre o gerador aleatório.} O enunciado original
da tarefa pedia o uso de \texttt{rand\_r()} (gerador thread-safe do
padrão POSIX) nas versões com contador compartilhado (V3 e V4).
Como o ambiente de desenvolvimento e testes utilizado é o Windows,
onde \texttt{rand\_r()} não está disponível de forma padronizada, essa
chamada foi substituída pela função \texttt{random\_double()}, baseada
no gerador Xorshift32 com estado (\texttt{seed}) privado por thread --
a mesma estratégia já adotada na Tarefa 8 para contornar a mesma
limitação. Do ponto de vista do experimento, essa troca não compromete
a análise: o objetivo da tarefa é comparar o \textbf{custo dos
mecanismos de sincronização} (critical, atomic, reduction), e não o
gerador de números aleatórios em si; usar o mesmo gerador
(Xorshift32) em V2, V3, V4 e V5 tem inclusive a vantagem de isolar
melhor essa variável, deixando a diferença de desempenho atribuível
quase exclusivamente à estratégia de sincronização escolhida.

Todas as versões foram executadas para o mesmo conjunto de tamanhos de
entrada ($n\_pontos \in \{10^4, 10^5, 10^6, 10^7, 10^8, 10^9\}$), com
20 threads, medindo o tempo de execução com \texttt{omp\_get\_wtime()}
e o erro absoluto/relativo em relação a $\pi$ real. O código-fonte
completo encontra-se em \texttt{../src/tarefa10.c} e é reproduzido no
Apêndice~A.

\section{Resultados}

A Tabela~\ref{tab:resultados} reproduz a saída obtida na execução do
programa (arquivo \texttt{output.txt}), com $\pi_{real} = 3.141592653589793$
e 20 threads.

\begin{table}[H]
\centering

\begin{tabular}{lccc}
\toprule
\textbf{N\_PONTOS} & \textbf{Versão} & \textbf{$\pi$ medido} & \textbf{Tempo (s)} \\
\midrule

10.000          & V1 & 3,11920000 & 0,0000 \\
10.000          & V2 & 3,14320000 & 0,0000 \\
10.000          & V3 & 3,14320000 & 0,0010 \\
10.000          & V4 & 3,14320000 & 0,0000 \\
10.000          & V5 & 3,14320000 & 0,0000 \\
\midrule

100.000         & V1 & 3,10040000 & 0,0010 \\
100.000         & V2 & 3,14408000 & 0,0000 \\
100.000         & V3 & 3,14408000 & 0,0040 \\
100.000         & V4 & 3,14408000 & 0,0030 \\
100.000         & V5 & 3,14408000 & 0,0000 \\
\midrule

1.000.000       & V1 & 3,13557200 & 0,0030 \\
1.000.000       & V2 & 3,14235200 & 0,0000 \\
1.000.000       & V3 & 3,14235200 & 0,0390 \\
1.000.000       & V4 & 3,14235200 & 0,0120 \\
1.000.000       & V5 & 3,14235200 & 0,0000 \\
\midrule

10.000.000      & V1 & 3,14204040 & 0,0090 \\
10.000.000      & V2 & 3,14160360 & 0,0020 \\
10.000.000      & V3 & 3,14160360 & 0,3640 \\
10.000.000      & V4 & 3,14160360 & 0,1210 \\
10.000.000      & V5 & 3,14160360 & 0,0020 \\
\midrule

100.000.000     & V1 & 3,14124624 & 0,0910 \\
100.000.000     & V2 & 3,14187084 & 0,0210 \\
100.000.000     & V3 & 3,14187084 & 3,7180 \\
100.000.000     & V4 & 3,14155864 & 1,1660 \\
100.000.000     & V5 & 3,14169476 & 0,0220 \\
\midrule

1.000.000.000   & V1 & 3,14137154 & 0,9300 \\
1.000.000.000   & V2 & 3,14161538 & 0,2110 \\
1.000.000.000   & V3 & 3,14161538 & 38,6540 \\
1.000.000.000   & V4 & 3,14157180 & 11,8630 \\
1.000.000.000   & V5 & 3,14160934 & 0,2110 \\

\bottomrule
\end{tabular}
\caption{Comparação de $\pi$ medido e tempo de execução entre as 5 versões, para 20 threads (extraído de \texttt{output.txt}).}
\label{tab:resultados}
\end{table}

Para facilitar a leitura, a Tabela~\ref{tab:speedup} isola o caso de
maior carga ($n\_pontos = 10^9$), tomando como referência a versão
mais lenta (V3, compartilhado + critical).

\begin{table}[H]
\centering
\begin{tabular}{lrr}
\toprule
\textbf{Versão} & \textbf{Tempo (s)} & \textbf{Speedup em relação a V3} \\
\midrule
V3 -- compartilhado + critical & 38,6540 & 1,0$\times$ (referência) \\
V4 -- compartilhado + atomic   & 11,8630 & $\approx$ 3,3$\times$ \\
V1 -- privado + critical (rand) & 0,9300  & $\approx$ 41,6$\times$ \\
V2 -- privado + critical (xorshift) & 0,2110 & $\approx$ 183,2$\times$ \\
V5 -- reduction                & 0,2110  & $\approx$ 183,2$\times$ \\
\bottomrule
\end{tabular}
\caption{Comparação de tempo para $n\_pontos = 1.000.000.000$.}
\label{tab:speedup}
\end{table}

\section{Análise}

\subsection{Comparação entre as cinco versões}

Os resultados da Tabela~\ref{tab:resultados} deixam claro que o erro
de estimativa (absoluto e relativo) é, como esperado, praticamente
idêntico entre V2, V3, V4 e V5 para um mesmo $n\_pontos$ -- afinal,
todas usam o mesmo gerador (Xorshift32) e a mesma quantidade de
pontos sorteados; pequenas diferenças de $\pi$ medido entre elas
(visíveis a partir de $n\_pontos=10^8$) vêm apenas da ordem de
acumulação das somas parciais em ponto flutuante/inteiro combinada
com a forma como as sementes de cada thread são geradas, e não da
qualidade do mecanismo de sincronização. V1 se destaca com erro maior
porque usa \texttt{rand()}, que tem baixa qualidade estatística e,
principalmente, não é thread-safe quando chamado concorrentemente
por várias threads sem proteção -- o que já era uma limitação
conhecida herdada da Tarefa 8.

O ponto central da tarefa, porém, é o \textbf{tempo de execução}, e
aí as diferenças são drásticas:

\begin{itemize}
    \item \textbf{V3 (compartilhado + critical) é, de longe, a pior
    versão.} Para $n\_pontos=10^9$ ela leva 38,65\,s, contra 0,21\,s
    de V2/V5 -- quase 200 vezes mais lenta. Isso acontece porque a
    seção crítica é disputada por até 20 threads \textbf{a cada
    acerto} (cerca de 78,5\% dos pontos sorteados caem dentro do
    círculo), e não apenas uma vez por thread. Cada entrada em
    \texttt{critical} envolve a aquisição de um mutex genérico, o
    que serializa fortemente as threads e faz o custo de
    sincronização dominar completamente o tempo total.

    \item \textbf{V4 (compartilhado + atomic) melhora bastante em
    relação a V3} (11,86\,s vs. 38,65\,s, um ganho de ~3,3$\times$),
    pois a diretiva \texttt{atomic} usa uma instrução atômica de
    hardware em vez de um
    mutex de propósito geral, com overhead bem menor para uma
    operação simples como \texttt{total\_acertos++}. Ainda assim, é
    muito mais lenta que as versões com contador privado, pois a
    variável continua sendo compartilhada e disputada a cada
    incremento, gerando contenção na mesma linha de cache
    (\emph{true sharing}) entre todas as threads.

    \item \textbf{V1 e V2 (contador privado + critical apenas ao
    final)} têm desempenho ordens de magnitude melhor que V3/V4,
    porque a seção crítica é executada apenas $O(\text{nº de
    threads})$ vezes (uma soma por thread ao final da região
    paralela), e não $O(n\_pontos)$ vezes. Isso confirma que o
    problema de V3/V4 não é o uso de \texttt{critical}/\texttt{atomic}
    em si, mas a \textbf{frequência} com que a sincronização é
    acionada.

    \item \textbf{V5 (reduction) tem desempenho equivalente ao melhor
    caso manual (V2)}, sem que o programador precise escrever
    nenhuma sincronização explícita. O compilador implementa a
    redução criando cópias privadas do contador por thread (evitando
    qualquer disputa durante o laço) e combinando-as ao final, de
    forma tipicamente mais eficiente ou pelo menos equivalente à
    melhor estratégia manual.
\end{itemize}

Em suma, a ordem de desempenho observada foi:
$$\text{V3 (critical/iteração)} \ll \text{V4 (atomic/iteração)} \ll \text{V1} \approx \text{V2} \approx \text{V5 (reduction)}.$$

\subsection{Reflexão: aplicabilidade em desempenho e produtividade}

Os mecanismos de sincronização do OpenMP não são substitutos
equivalentes entre si -- cada um representa um compromisso diferente
entre \textbf{desempenho}, \textbf{generalidade} e
\textbf{produtividade/simplicidade do código}:

\begin{itemize}
    \item \texttt{reduction} é o mecanismo de \textbf{maior
    produtividade e, para o caso testado, também de melhor
    desempenho}: o programador não escreve nenhum código de
    sincronização, elimina uma classe inteira de bugs (condições de
    corrida, esquecimento de proteger uma variável, deadlocks) e
    delega ao compilador/runtime a escolha da melhor estratégia de
    combinação. A limitação é o escopo: só se aplica a operações de
    redução suportadas pela cláusula (soma, produto, mínimo, máximo,
    operadores lógicos/bit a bit sobre variáveis escalares), não
    servendo para proteger estruturas de dados complexas ou lógica
    condicional arbitrária.

    \item \texttt{atomic} é indicado quando a operação protegida é
    \textbf{uma única instrução simples} sobre uma variável escalar
    (incremento, soma, subtração, min/max, etc.). Tem
    overhead menor que \texttt{critical} porque usa instruções
    atômicas de hardware, mas ainda paga o preço da disputa pela
    mesma linha de cache quando muitas threads o atualizam
    simultaneamente com alta frequência (como ficou evidente em V4).

    \item \texttt{critical} (anônima) é o mecanismo mais
    \textbf{genérico e de maior overhead}: protege blocos de código
    arbitrários, com múltiplas instruções, variáveis ou até chamadas
    de função, mas paga o custo de um mutex de propósito geral e,
    pior, \textbf{todas} as regiões críticas anônimas do programa
    compartilham implicitamente o mesmo lock global -- ou seja, uma
    seção crítica não relacionada a outra ainda pode bloqueá-la
    desnecessariamente. Deve ser reservado para blocos que realmente
    não podem ser expressos como \texttt{atomic}/\texttt{reduction}, e
    sempre \textbf{minimizando a frequência de entrada} (idealmente
    acumulando em variáveis privadas e entrando na seção crítica uma
    única vez por thread, como nas versões V1/V2, e nunca dentro do
    laço principal, como em V3).

    \item \texttt{critical} \textbf{nomeada}
    (\texttt{\#pragma omp critical(nome)}) resolve exatamente a
    limitação do lock global único: quando existem \textbf{duas ou
    mais seções críticas protegendo dados independentes} (por
    exemplo, um contador de acertos e, separadamente, um log de
    erros ou um acumulador de estatísticas), nomear cada seção
    crítica permite que threads em seções com nomes diferentes
    executem \textbf{concorrentemente}, em vez de serem serializadas
    por um lock global que não tem relação com o dado que estão
    protegendo. Deve ser usada sempre que houver múltiplas regiões
    críticas logicamente independentes coexistindo no mesmo programa.

    \item \textbf{Locks explícitos} são a opção de \textbf{menor
    produtividade, porém maior controle}: úteis quando a granularidade
    da proteção precisa ser definida dinamicamente em tempo de
    execução (por exemplo, um vetor de locks, um por posição/bucket
    de uma tabela hash, permitindo que threads acessando posições
    diferentes não se bloqueiem entre si), quando é necessário tentar
    adquirir a proteção sem bloquear (\texttt{omp\_test\_lock}), ou
    quando o lock precisa ser mantido por mais tempo do que o
    permitido de forma limpa por um bloco \texttt{critical}/\texttt{atomic}
    (por exemplo, atravessando chamadas de função). O custo é a
    responsabilidade adicional do programador: inicializar e destruir
    corretamente cada lock, e evitar deadlocks ao usar múltiplos locks
    simultaneamente.
\end{itemize}

\subsection{Roteiro proposto para escolha do mecanismo}

Com base nos resultados obtidos e na discussão acima, propõe-se o
seguinte roteiro de decisão, a ser percorrido nesta ordem:

\begin{enumerate}
    \item \textbf{A operação é uma redução simples sobre uma
    variável escalar} (soma, produto, mínimo, máximo, and/or/xor)?
    \\ $\rightarrow$ use \texttt{reduction}. É a opção mais rápida e
    mais segura, e deve ser sempre a primeira escolha quando aplicável.

    \item \textbf{A operação protegida envolve múltiplas variáveis,
    uma estrutura de dados, ou lógica condicional não redutível a uma
    única instrução} (por exemplo, atualizar uma \texttt{struct}, uma
    lista encadeada ou um vetor de tamanho variável)? \\
    $\rightarrow$ use \texttt{critical}, mas \textbf{minimize sua
    frequência de entrada}: sempre que possível, acumule os
    resultados em uma variável/estrutura privada por thread durante o
    laço e entre na seção crítica apenas uma vez, ao final da região
    paralela (como em V1/V2), nunca a cada iteração (como em V3).

    \item \textbf{Existem duas ou mais seções críticas protegendo
    dados independentes no mesmo programa}? \\
    $\rightarrow$ nomeie cada uma com
    \texttt{\#pragma omp critical(nome)}, para que a disputa por uma
    delas não bloqueie threads que precisam apenas da outra.

    \item \textbf{A operação é uma única instrução simples sobre um
    único escalar} (incremento, soma, min/max) e, por algum motivo,
    não se encaixa em \texttt{reduction} (ex.: a variável também é
    lida/usada de outras formas dentro da região paralela)? \\
    $\rightarrow$ prefira \texttt{atomic} a \texttt{critical}: mesmo
    semântica de exclusão mútua, porém com overhead bem menor por
    usar instruções atômicas de hardware.

    \item \textbf{É necessário controle fino sobre a proteção}? \\
    $\rightarrow$ use \textbf{locks explícitos}
    (\texttt{omp\_lock\_t}), aceitando o custo adicional de
    gerenciá-los corretamente em troca da flexibilidade que nenhum
    dos mecanismos anteriores oferece.
\end{enumerate}

De forma resumida: a \textbf{produtividade} decresce e a
\textbf{complexidade de uso} cresce na ordem
$$\texttt{reduction} \;\rightarrow\; \texttt{critical (uso disciplinado)} \;\rightarrow\;$$
$$\texttt{atomic} \;\rightarrow\; \texttt{critical nomeada} \;\rightarrow\; \texttt{locks explícitos}$$

Enquanto o \textbf{desempenho}, para o caso de agregar um contador
escalar como o deste experimento, tende a seguir a ordem inversa até
certo ponto: \texttt{reduction} e contador privado com
\texttt{critical} ao final empatam como melhores opções, seguidos por
\texttt{atomic} por iteração e, por último, \texttt{critical} por
iteração -- confirmado empiricamente pela diferença de quase 200
vezes entre V3 e V5 para $n\_pontos = 10^9$.

\section{Conclusão}

O experimento confirmou, de forma quantitativa, uma regra prática
importante em programação paralela com OpenMP: \textbf{o mecanismo de
sincronização escolhido importa menos do que a frequência com que ele
é acionado}. Um \texttt{critical} usado uma vez por thread (V1/V2) é
ordens de magnitude mais rápido que o mesmo \texttt{critical} usado a
cada iteração de um laço com centenas de milhões de iterações (V3),
mesmo protegendo exatamente a mesma informação. Trocar
\texttt{critical} por \texttt{atomic} reduz o dano quando a
sincronização por iteração é inevitável, mas não elimina a contenção
gerada por uma variável compartilhada disputada por muitas threads.

A cláusula \texttt{reduction}, quando aplicável, é a melhor escolha
tanto em desempenho quanto em produtividade, pois evita completamente
a necessidade de o programador implementar manualmente a
acumulação privada e a combinação final -- e, neste experimento,
igualou o desempenho da melhor implementação manual (V2) sem exigir
nenhuma linha de código de sincronização. Mecanismos mais
flexíveis, como \texttt{critical} nomeada e locks explícitos, seguem
sendo ferramentas importantes para os casos em que \texttt{reduction}
e \texttt{atomic} não se aplicam, mas devem ser reservados para essas
situações específicas, e não usados por padrão quando alternativas
mais simples e eficientes estão disponíveis.

\newpage

\appendix
\section{Código-Fonte}

O código completo utilizado nesta tarefa está localizado em
\texttt{../src/tarefa10.c} e é reproduzido a seguir.

\lstinputlisting[style=cstyle, caption={Código-fonte completo -- \texttt{tarefa10.c}}, label={lst:tarefa10}]{../src/tarefa10.c}

\end{document}